% problem-000366 1 晶体结构与晶胞
\section*{1. 晶体结构与晶胞}
\textit{下列为分问。}

\subsection*{1-1}
钒酸盐与磷酸盐结构相似。请画出 $\mathrm { V O } _ { 4 } \mathrm { \dot { \Omega } }$ 3−、 $\mathrm { H _ { 2 } V O _ { 4 } } ^ { - }$ $\mathrm { V O _ { 2 } ( H _ { 2 } O ) _ { 4 } } ^ { + }$ ，和 $\mathrm { V _ { 2 } O _ { 7 } } ^ { 4 - }$ −的空间构型。

\subsection*{1-2}
⽣理条件下的钒以多种氧化态存在，各种氧化态可以相互转化。通常细胞外的钒是 $\mathrm { V } ( + \mathrm { V } )$ ，⽽细胞内的钒是 $\mathrm { V } ( + \mathrm { I V } )$ 。研究表明，钒酸⼆氢根离⼦可与亚铁⾎红素 $\mathrm { ( M t r c - F e ^ { 2 + } ) }$ 反应，写出该反应的离⼦⽅程式。

\subsection*{1-3}
-1 已知配合物 $\mathrm { \small { [ V O N ( C H _ { 2 } C O O ) _ { 3 } ] } }$ 在⽔溶液中的⼏何构型是唯⼀的，画出它的空间构型图。

\subsection*{1-3}
-2 理论推测上述配合物分⼦在晶体中是有⼿性的，指出产⽣⼿性的原因。

\subsection*{1-4}
钒酸钇晶体是近年来新开发出的优良双折射光学晶体，在光电产业中得到⼴泛应⽤。可以在弱碱性溶液中⽤偏钒酸铵和硝酸钇合成。写出以 $\mathrm { Y _ { 2 } O _ { 3 } }$ 与 $\mathrm { { V } _ { 2 } O _ { 3 } }$ 为主要原料合成钒酸钇的化学⽅程式。

\subsection*{1-5}
若以市售分析纯偏钒酸铵为原料制备⾼纯钒酸钇单晶，需将杂质铁离⼦含量降⾄⼀定数量级。设每升偏钒酸铵溶液中含三价铁离⼦为 $5 . 0 { \times } 1 0 ^ { - 5 } \mathrm { m o l }$ ，⽤0.01 mol L−1的鏊合剂除铁。

\subsection*{1-5}
-1 说明不采取使铁离⼦⽔解析出沉淀的⽅法除铁的理由。

\subsection*{1-5}
-2 通过计算说明如何选择螯合剂使偏钒酸铵含铁量降⾄10−30 mol L−1以下。

\textit{（原卷此处含表格/仪器试剂表，详见 PDF 或 MD 源文件。）}

\textit{（原卷此处含表格/仪器试剂表，详见 PDF 或 MD 源文件。）}
